% Workaround para criação de capítulo não-numerado alinhado à margem
\chapter*{}
\noindent
\phantomsection{\MakeUppercase{\textbf{Introdução}}}
\addcontentsline{toc}{chapter}{INTRODUÇÃO}

% \newline
% \newline

%\section{Motivação}

\vspace{1.4cm}

Os setores da agropecuária e da mineração têm sido pilares fundamentais na história econômica do Brasil, contribuindo para o desenvolvimento tanto no âmbito regional quanto nacional. Em 2021, com as commodities em alta em meio à pandemia, a soma de agro e mineração superou a de manufatura no PIB brasileiro pela primeira vez em décadas e a tendência se manteve em 2022 com os efeitos da guerra da Ucrânia \cite{BBC2023}.

Sendo assim, com a crescente demanda por recursos minerais essenciais, como argila, calcário e ferro, a mineração não apenas fornece matérias-primas para diversas indústrias, mas também gera empregos e contribui para a economia do Brasil. Com isso, o setor de mineração volta ao foco como uma das atividades econômicas importantes do Brasil, contribuindo significativamente para o PIB nacional e para a geração de empregos \cite{ibram2020tecnologias}.

No entanto, a atividade mineradora também traz consigo desafios e riscos significativos, especialmente no que diz respeito ao gerenciamento de resíduos e de rejeitos resultantes dos processos de beneficiamento do minério, devido aos seus impactos ambientais e sociais \cite{ipea2017boletim}.

A gestão de rejeitos de mineração é uma questão crítica que requer abordagens inovadoras e tecnologias avançadas para garantir a segurança e a sustentabilidade ambiental das operações \cite{ibram2020tecnologias}. Pois a má gestão dos rejeitos pode levar a uma série de consequências adversas, incluindo mudanças ambientais, desvalorização de imóveis e até mesmo desastres como os que ocorreram em Mariana e Brumadinho \cite{ipea2017boletim}, que serviram como um chamado de alerta para a necessidade urgente de melhorar os sistemas de segurança e prevenção em barragens de rejeitos.

Nesse contexto, a engenharia da computação oferece uma variedade de soluções inovadoras que podem ser aplicadas para aumentar a segurança e reduzir os riscos associados às barragens de rejeitos. A prevenção de desastres em barragens de rejeitos requer a implementação de tecnologias modernas e sistemas de monitoramento sofisticados para identificar e mitigar os riscos potenciais \cite{icmm2020global}.

Como exemplo de aplicação, este trabalho utilizou um higrômetro de resistência para solo, implementado na plataforma \textit{open-source} Arduino, como instrumentação geotécnica de baixo custo para o monitoramento contínuo da umidade do solo, integrado a dados meteorológicos do Centro Nacional de Monitoramento e Alertas de Desastres Naturais (CEMADEN) e do Instituto Nacional de Meteorologia (INMET) para correlacionar a umidade detectada com parâmetros climatológicos como precipitação acumulada e previsão pluviométrica, contribuindo para a segurança das operações de mineração por meio da emissão automatizada de alertas graduais de risco, com potencial de ajudar a prevenir desastres semelhantes aos de Mariana e Brumadinho.

O Capítulo 1 apresenta a revisão bibliográfica sobre a atividade de mineração, os tipos de barragens de rejeitos, a legislação vigente e os acidentes de Mariana e Brumadinho. O Capítulo 2 traz a fundamentação teórica sobre instrumentação geotécnica e automação de sistemas de monitoramento. O Capítulo 3 apresenta o objetivo geral e os objetivos específicos do trabalho. O Capítulo 4 descreve os materiais e métodos utilizados no desenvolvimento do sistema. O Capítulo 5 apresenta os resultados obtidos e o Capítulo 6 traz a discussão desses resultados, contemplando a cobertura dos objetivos específicos e a discussão dos resultados. Por fim, são apresentadas as conclusões, seguidas das referências utilizadas no trabalho.

\vspace{10mm} %5mm vertical space